\section{First-Order Phase Transitions and Phase Diagrams}

\paragraph{Phases and phase transitions.}
A \textcolor{blue}{phase} of a system is a region of the space of natural
variables in which the thermodynamic potential is analytic, that is,
infinitely differentiable, with no kinks or divergences. A
\textcolor{blue}{phase transition} is a point where the potential fails to
be analytic. Strictly, such singularities appear only in the thermodynamic
limit $N\to\infty$, because the partition function of a finite system is
analytic.

\paragraph{The Gibbs free energy and the composition variables.}
Phases are usually studied at fixed particle numbers, pressure and
temperature for two reasons. Experimentally, these variables are easy to
control. Structurally, the phases inside such a system exchange energy,
volume and particles with each other, so $(\{\mu_i\},P,T)$ are equal in all
phases at equilibrium. The natural choice of potential is then the Gibbs
free energy $G(T,P,\{N_i\})$. Because $G$ is extensive,
\begin{align*}
    G(T,P,\{N_i\}) = N\,g(T,P,\{x_i\}), \qquad x_i = N_i/N,
\end{align*}
so the total size $N$ does not affect the analyticity of $G$, and only $g$
matters. Of the $C$ mole fractions $\{x_i\}$ of a $C$-component system, only
$C-1$ are independent, because $\sum_i x_i=1$. A phase is therefore a region
of the intensive variables $(T,P,\{x_i\})$, which span a
$(C-1)+2=(C+1)$-dimensional space. The collection of all phases in this
space is the \textcolor{blue}{phase diagram}.

\paragraph{First-order transitions as creases of $g$.}
At equilibrium, $g$ is continuous, so a phase transition is never a jump in
the potential itself. With $g$ as one more axis on the $(C+1)$-dimensional
space $(T,P,\{x_i\})$, the potential forms a hypersurface
$g(T,P,\{x_i\})$. A \textcolor{blue}{first-order transition} is a crease
(kink) of this hypersurface, across which $g$ is continuous but the first
derivatives of $g$ jump:
\begin{align*}
    s = -\frac{\partial g}{\partial T}, \qquad v = \frac{\partial g}{\partial P},
\end{align*}
which are the entropy and the volume per particle. Equivalently,
$S=-\partial G/\partial T$ and $V=\partial G/\partial P$ are discontinuous,
which gives a latent heat $L = T\Delta S$ and a volume change $\Delta V$.
Melting and boiling are examples. On the transition locus, the coexisting
phases share the same $(\{\mu_i\},P,T)$ but differ in $s$, $v$ and
composition. For this reason, the transition appears as the boundary of a
coexistence region in the phase diagram.

\paragraph{Higher-order transitions.}
Higher-order transitions also exist. The order of a transition is the order
of the lowest derivative of $g$ that is discontinuous or singular. For
higher-order transitions, the hypersurface has no kink, and the singularity
appears in the curvature. These transitions include the continuous
transitions, such as the ferromagnetic Curie point, where response functions
such as the heat capacity diverge. A later section treats these
transitions.

\subsection{Gibbs Phase Rule}

\paragraph{Counting the degrees of freedom.}
Suppose $\mathcal{P}$ phases of a system with $C$ components coexist. At
equilibrium, all the phases share the same $C+2$ intensive variables
$(T,P,\mu_1,\dots,\mu_C)$. Each phase has its own Gibbs--Duhem relation,
\begin{align*}
    S_\alpha\,dT - V_\alpha\,dP + \sum_i N_{i,\alpha}\,d\mu_i = 0,
    \qquad \alpha = 1,\dots,\mathcal{P},
\end{align*}
which imposes $\mathcal{P}$ independent constraints on these variables. The
number of degrees of freedom $F$ is the number of variables minus the number
of constraints,
\begin{align}
    \boxed{F+\mathcal{P}=C+2.}
\end{align}
This relation is the \textcolor{blue}{Gibbs phase rule}. The rule counts
only intensive degrees of freedom, so the amount of each phase remains free
and is not counted.

\paragraph{Limits of the phase rule.}
The rule in this form applies only if $(T,P,\mu_1,\dots,\mu_C)$ are the only
intensive variables. Magnetic systems, for example, also depend on the
magnetic field $\vec B$, which adds another intensive variable to the count.

\subsection{Clapeyron Equation and Fixed-Composition ($P$--$T$) Phase Diagram}

\paragraph{The two-phase coexistence curve.}
Temperature and pressure appear in every phase diagram, whatever the number
of components. The $P$--$T$ plane at fixed composition $\{x_i\}$ is a
two-dimensional slice of the $(C+1)$-dimensional phase diagram. Only $P$
and $T$ vary in this slice, so the slice has two degrees of freedom. Two
phases, $1$ and $2$, coexist when every chemical potential matches,
\begin{equation*}
    \mu_i^{(1)}(T,P) = \mu_i^{(2)}(T,P), \qquad i=1,\dots,C.
\end{equation*}
These conditions relate $P$ and $T$ and reduce the degrees of freedom from
two to one. The coexistence locus is therefore a curve in the plane.

\paragraph{The Clapeyron equation.}
Along this curve, the two phases keep equal chemical potentials, so
$d\mu_i^{(1)}=d\mu_i^{(2)}$, and the Gibbs--Duhem relation at fixed
composition gives
\begin{equation*}
    -s_1\,dT + v_1\,dP = -s_2\,dT + v_2\,dP ,
\end{equation*}
which is the \textcolor{blue}{Clapeyron equation}
\begin{equation}
    \boxed{\frac{dP}{dT}=\frac{\Delta s}{\Delta v}}, \qquad
    \Delta s = s_2-s_1,\;\; \Delta v = v_2-v_1 .
\end{equation}
At coexistence, the transition proceeds reversibly at fixed $T$ and $P$, so
the latent heat per particle is $\ell = T\Delta s$, and
\begin{equation}
    \boxed{\frac{dP}{dT}=\frac{\ell}{T\,\Delta v}.}
\end{equation}

\paragraph{Three-phase coexistence and the triple point.}
If a third phase, $3$, also coexists, then
\begin{equation*}
    \mu_i^{(1)}(T,P)=\mu_i^{(2)}(T,P)=\mu_i^{(3)}(T,P), \qquad i=1,\dots,C .
\end{equation*}
Each pair of phases has its own coexistence curve, $\mu^{(1)}=\mu^{(2)}$,
$\mu^{(2)}=\mu^{(3)}$ and $\mu^{(1)}=\mu^{(3)}$. A point on the first two
curves satisfies $\mu^{(1)}=\mu^{(2)}=\mu^{(3)}$, so the point lies on the
third curve as well, and only two of the three conditions are independent.
Two independent equations in the two unknowns $(T,P)$ leave no degrees of
freedom, so three phases can coexist only at an isolated point, the
\textcolor{blue}{triple point}, where the three two-phase curves intersect.

\subsection{Unary System}

\paragraph{Definition and phase rule.}
A \textcolor{blue}{unary system} has only one component, $C=1$, and is a
special case of the fixed-composition slice with $x=1$. The unary system is
the building block of higher-dimensional phase diagrams and introduces the
terminology used for them. For $C=1$, the phase rule reads
$F = 3-\mathcal{P}$, so the number of coexisting phases $\mathcal{P}$ fixes
the number of free variables.

\paragraph{Bivariant, univariant and invariant equilibria.}
Within a single phase ($\mathcal{P}=1$, $F=2$), $P$ and $T$ vary freely,
and the system is in \textcolor{blue}{bivariant equilibrium}. When two
phases coexist ($\mathcal{P}=2$, $F=1$), only one of $P$ and $T$ varies
freely, and the coexistence curve $P(T)$, with the slope given by the
Clapeyron equation, fixes the other. The system is then in
\textcolor{blue}{univariant equilibrium}, and the two-phase coexistence
locus is a \textcolor{blue}{univariant curve}. When three phases coexist
($\mathcal{P}=3$, $F=0$), no variable varies freely. The system is in
\textcolor{blue}{invariant equilibrium}, and the point of three-phase
coexistence is an \textcolor{blue}{invariant point}, the triple point.
